% !TEX program = xelatex
% 2026 年高教社杯全国大学生数学建模竞赛 A 题
% 摘要、问题重述、问题分析、模型假设与符号、模型评价与推广

\documentclass[12pt,a4paper]{ctexart}

\usepackage{amsmath,amssymb,bm}
\usepackage{booktabs}
\usepackage{multirow}
\usepackage{array}
\usepackage{geometry}
\usepackage{caption}

\geometry{left=2.7cm,right=2.7cm,top=2.6cm,bottom=2.6cm}
\captionsetup{font=small,labelsep=quad}
\linespread{1.32}

\begin{document}

\begin{center}
  {\zihao{3}\bfseries 药材烘干过程的传热传质建模与数值仿真}
\end{center}

\vspace{0.5em}
\begin{center}
  {\zihao{4}\bfseries 摘\quad 要}
\end{center}

\noindent
中药材烘干是决定成品品质的关键工序。热风烘干包含预热平衡与恒温干燥两个阶段，
工艺参数选取不当会造成干燥效率低、能耗高、品质不稳定。本文针对圆柱形药材，
建立了径向耦合传热传质数学模型，并对四个问题给出定量解答。

\textbf{针对问题一}，考虑药材长径比达 6.25、端面仅占侧面 8\%，
将其简化为无限长圆柱的一维径向轴对称问题。由附录 2 的物性算得热扩散率
$\alpha=1.689\times10^{-7}$\,m$^2$/s 与水分扩散系数
$D(C_0)=4.938\times10^{-9}$\,m$^2$/s 相差 34 倍，
1800\,s 内热扩散长度 17.4\,mm 已接近半径而水分扩散长度仅 2.98\,mm，
由此判断该阶段以达热平衡为主。关键的结构性质是：$D$ 只依赖 $C$ 不含 $T$、
其余物性为常数，故\textbf{热方程与水分方程完全解耦}，可分别独立求解。
采用节点中心有限体积法配合后向 Euler 与 Picard 迭代求解，水分守恒残差为
$6.8\times10^{-14}$（机器精度）。结果表明：1800\,s 内药材中心含水率始终保持
2.5500\,kg/kg，而表面已降至 1.5103，水分变化局限在最外侧约 $3\sim5$\,mm；
温度则在整个截面上重新分布，中心 33.5765\,$^\circ$C、表面 36.7863\,$^\circ$C，
径向温差 3.21\,$^\circ$C。

\textbf{针对问题二}，附录 3 给出 $\rho$、$c_p$、$k$ 随 $C$ 变化、
$D$ 随 $C$ 与 $T$ 变化的经验式，两方程由解耦转为强耦合，必须联立迭代。
由附件 1 末段 1\,h 的平均值（49.9989\,$^\circ$C、0.04999\,kg/kg）
确定恒温干燥段工况为 50\,$^\circ$C 与 0.05\,kg/kg。结果表明药材在约 2.5\,h
内基本达到烘房温度（3\,h 时径向温差仅 0.12\,$^\circ$C），含水率则依
"由表及里"的次序下降：0.5\,h 时中心基本未变而表面已降至 1.6491，
3\,h 时中心 1.7663、表面 1.0082，全过程始终由内部扩散控制。

\textbf{针对问题三}，烘干判据"各处含水率低于 0.15\,kg/kg"等价于中心处达标，
即 $t_{\mathrm{dry}}=\min\{t\mid C(0,t)<0.15\}$。求得
\textbf{烘干时长 57.42\,h（2.392 天）}，落在题面所述 $2\sim3$ 天区间内；
其中 36\,h 之后中心由 0.1858 缓慢逼近 0.15 竟耗时约 21\,h，
占总时长的 37\%，是工程上最不经济的阶段。

\textbf{针对问题四}，附件 2 显示半径由 2.000\,cm 收缩至 1.198\,cm，
且收缩高度前置（50\% 的收缩在 2.5\,h 内完成，99\% 在 22\,h 内完成），
而含水率的缓慢下降持续到 50\,h 之后，说明二者不存在简单对应关系，
故直接采用实测 $R(t)$。引入贴体坐标 $\xi=r/R(t)$ 将动边界化为固定区域，
方程相应增加对流项 $\xi R\dot R\,\partial C/\partial\xi$；
$R(t)$ 用保单调的三次 Hermite 插值以避免 $\dot R$ 跳变。
求得\textbf{烘干时长 47.73\,h（1.989 天）}，比不考虑收缩时缩短 16.9\%；
作为对照，同一物性下固定半径的烘干时长为 5.41 天，可见
\textbf{忽视收缩会将烘干时长高估约 1.7 倍}。

\textbf{本文的特色}在于：一是严格区分了问题一中热质解耦与问题二中热质耦合的
结构性差异；二是对题面未给出的闭合关系（传质边界条件的密度口径、汽化潜热）
作了定量讨论，并依据湿球温度下限判断出题面隐含忽略汽化潜热；
三是全部结果通过了守恒性、网格与时间步收敛性以及独立求解器的交叉验证。
四个问题算得的烘干时长均与题面给出的 2--3 天量级自洽。

\vspace{1em}
\noindent\textbf{关键词}：\ 热风干燥；耦合传热传质；圆柱坐标；有限体积法；
动边界；收缩；汽化潜热

\newpage

\section{问题重述}

\subsection{问题背景}

干燥是决定中药材成品品质的关键工序之一，热风烘干是常见方式，
包括预热平衡与恒温干燥两个阶段。工艺参数选取不当会导致干燥效率低、
能耗高、成品品质不稳定；传统试验优化成本高、周期长，
故需借助数理分析与数值仿真获得干燥规律。

\subsection{问题提出}

某圆柱形中药材长 25\,cm、半径 2\,cm，初始温度 28\,$^\circ$C、
初始干基含水率 2.55\,kg/kg。烘房空气的温度与水分浓度见附件 1，
附录 2 给出问题一的物性参数，附录 3、4 分别给出问题二/三与问题四的
经验公式，附件 2 给出烘干过程中药材半径的变化。需要解决：

\begin{enumerate}
  \item 建立预热平衡阶段药材温度与水分浓度变化的数学模型，
        给出 1800\,s 内按指定时空格点采样的结果；
  \item 建立整个烘干过程（预热平衡 $+$ 恒温干燥）的数学模型，
        给出 3\,h 内的结果；
  \item 按"各处含水率低于 0.15\,kg/kg"的要求确定烘干所需时间；
  \item 考虑药材因失水发生的尺寸变化，重新确定烘干时长。
\end{enumerate}

\section{问题分析}

四个问题本质上求解同一个物理过程的不同层面：

\begin{table}[htbp]
  \centering
  \caption{四个问题的差异}
  \label{tab:overview}
  \small
  \begin{tabular}{llll}
    \toprule
    问题 & 物性来源 & 时间跨度 & 结构性特点 \\
    \midrule
    一 & 附录 2（恒定）        & 30 min  & 热、质方程解耦 \\
    二 & 附录 3（随 $C$、$T$） & 2--3 天 & 强耦合，需联立迭代 \\
    三 & 附录 3                & 2--3 天 & 反求达标时刻 \\
    四 & 附录 4（随 $C$、$T$） & 2--3 天 & 动边界，需坐标变换 \\
    \bottomrule
  \end{tabular}
\end{table}

问题的物理内核是\emph{含湿多孔介质内的耦合传热传质}：热量由表面向内部传导，
水分由内部向表面迁移并在表面被流动的空气带走，两者通过随含水率、
温度变化的物性相互影响。

三个量级分析决定了建模方向。其一，长径比与端面占比说明可退化为
一维径向问题；其二，热扩散率与水分扩散系数相差数十倍，
说明温度与水分的时间尺度不同，预热阶段以达热平衡为主；其三，
两个 Biot 数均处于中值区间，说明不能作集总参数近似，
必须给出沿半径的分布。

\section{模型假设}

\begin{enumerate}
  \item[(A1)] 药材均质、各向同性，孔隙水与固体骨架局部热平衡；
  \item[(A2)] 忽略轴向输运，问题简化为一维径向轴对称；
  \item[(A3)] 无内热源、忽略辐射，仅考虑表面对流换热与对流传质；
  \item[(A4)] 问题一至三不考虑尺寸变化，问题四按附件 2 的实测半径收缩；
  \item[(A5)] 物性按各问对应附录取值；问题一中 $\rho$、$c_p$、$k$ 为常数
        且 $D=D(C)$，问题二起 $D=D(C,T)$、$\rho,c_p,k$ 随 $C$ 变化；
  \item[(A6)] 对流换热系数与对流传质系数题面未给出，统一沿用附录 2 的
        $h=25$\,W/(m$^2\cdot$K)、$h_m=8\times10^{-7}$\,m/s；
  \item[(A7)] 忽略汽化潜热对能量平衡的贡献。若按同一密度口径计入，
        表面温度会被压至 $6.1\,^\circ$C，低于由空气状态算得的湿球温度
        $34.7\,^\circ$C，与热力学第二定律矛盾，故判定题面所给参数体系
        隐含该项可略；
  \item[(A8)] 水分的输运用 Fick 型扩散描述，表面传质用第三类边界条件描述。
\end{enumerate}

\section{符号说明}

\begin{table}[htbp]
  \centering
  \caption*{主要符号}
  \small
  \begin{tabular}{cll}
    \toprule
    符号 & 含义 & 单位 \\
    \midrule
    $r,\ t$            & 径向坐标、时间              & m, s \\
    $T,\ C$            & 药材温度、干基含水率        & $^\circ$C, kg/kg \\
    $R,\ L$            & 药材半径、长度              & m \\
    $\xi$              & 贴体坐标 $\xi=r/R(t)$       & --- \\
    $T_\infty,C_\infty$ & 烘房空气温度、水分浓度     & $^\circ$C, kg/kg \\
    $\rho,\ c_p,\ k$   & 密度、比热容、热传导系数    & kg/m$^3$, J/(kg$\cdot$K), W/(m$\cdot$K) \\
    $D$                & 水分扩散系数                & m$^2$/s \\
    $h,\ h_m$          & 对流换热、传质系数          & W/(m$^2\cdot$K), m/s \\
    $\alpha$           & 热扩散率                    & m$^2$/s \\
    $Bi,\ Bi_m$        & 传热、传质 Biot 数          & --- \\
    $t_{\mathrm{dry}}$ & 烘干时长                    & s \\
    \bottomrule
  \end{tabular}
\end{table}

\section{模型的评价与推广}

\subsection{模型优点}

\begin{enumerate}
  \item 从控制方程出发刻画传热传质机理，而非经验拟合，
        参数全部来自题面，物理意义明确；
  \item 明确识别出问题一中热质解耦这一结构性质，
        使该问的计算量降为问题二的三分之一；
  \item 采用有限体积离散，格式本身严格守恒，
        水分守恒残差达 $6.8\times10^{-14}$；
  \item 用贴体坐标处理收缩引起的动边界，把变区域问题化为固定区域问题，
        无需网格重划；
  \item 对题面未给出的闭合关系逐一作了定量辨析，
        并用湿球温度下限否决了含汽化潜热的方案，
        避免了物理上不可能的结果。
\end{enumerate}

\subsection{模型缺点}

\begin{enumerate}
  \item 假设药材均质、各向同性，未考虑真实的孔隙结构与各向异性；
  \item 将水分输运简化为单一 Fick 扩散，未区分自由水与结合水的差异，
        在含水率很低时可能与实际有偏差；
  \item 附件 1 只覆盖 4\,h，恒温段取常数为假设，
        未考虑烘房工况可能存在的漂移；
  \item 传质边界条件的密度口径题面未明确，不同口径下表面含水率相差约 40\%，
        该不确定性无法由题面数据消除。
\end{enumerate}

\subsection{模型推广}

\begin{enumerate}
  \item \textbf{变工况优化}：把烘房温度与风速作为控制变量，
        以能耗与品质为目标建立优化模型，寻找分段变温的干燥曲线，
        减少问题三中指出的"最后 25 小时只换来很小含水率下降"的低效段；
  \item \textbf{多维扩展}：对长径比较小的药材，可将模型扩展为二维
        轴对称或三维问题，只需更换相应的散度算子；
  \item \textbf{模型参数反演}：以实测含水率分布为目标，
        用贝叶斯或最小二乘方法反演 $h$、$h_m$ 与 $D$ 的经验常数，
        可提高对其他药材品种的适用性；
  \item \textbf{方法推广}：本文"贴体坐标 $+$ 保单调插值 $+$
        有限体积隐式推进"的处理框架同样适用于其他带收缩或膨胀的
        传热传质过程，如食品干燥、污泥干化、锂电池电极干燥等。
\end{enumerate}

\section*{参考文献}
\addcontentsline{toc}{section}{参考文献}

\begin{enumerate}
  \item[{[1]}] 王补宣. 工程传热传质学[M]. 北京: 科学出版社, 1998.
  \item[{[2]}] 潘永康, 王喜忠, 刘相东. 现代干燥技术[M]. 2 版. 北京: 化学工业出版社, 2007.
  \item[{[3]}] 陶文铨. 数值传热学[M]. 2 版. 西安: 西安交通大学出版社, 2001.
  \item[{[4]}] Patankar S V. Numerical Heat Transfer and Fluid Flow[M].
        New York: Hemisphere Publishing, 1980.
  \item[{[5]}] Fritsch F N, Carlson R E. Monotone piecewise cubic interpolation[J].
        SIAM Journal on Numerical Analysis, 1980, 17(2): 238--246.
  \item[{[6]}] Crank J. The Mathematics of Diffusion[M]. 2nd ed.
        Oxford: Clarendon Press, 1975.
\end{enumerate}

\end{document}
